\section{问题提出的背景}


\subsection{背景介绍}

近年来，现代处理器为了追求极致的计算性能，广泛采用了预测执行、乱序执行以及多级缓存等复杂的微架构设计。
然而，这些底层性能优化技术在提升指令吞吐量的同时，也打破了传统软件层面的安全隔离边界，引发了一系列隐蔽且危害极大的处理器微架构侧信道漏洞。
自2018年 Meltdown\cite{10.1145/3357033} 与 Spectre\cite{8835233} 漏洞被首次披露以来，硬件级侧信道攻击迅速成为系统安全领域的重大威胁。
攻击者无需获取高级系统特权，即可通过观测缓存访问延迟、分支预测器状态等微架构物理特征，推导出程序执行过程中的敏感信息流，甚至能够读取一台物理机器上所有地址空间的内存数据。
随后数年间，诸如Foreshadow\cite{10.5555/3277203.3277277}、ZombieLoad\cite{10.1145/3319535.3354252}、
Retbleed\cite{281436} 以及近期的 Downfall\cite{10.5555/3620237.3620639} 和 Inception\cite{10.5555/3620237.3620646} 等新型变种攻击层出不穷。
这些攻击不仅波及了Intel、AMD、ARM、RISCV 等主流指令集架构的各类处理器，更直接威胁到了云计算隔离、可信执行环境以及现代操作系统的核心安全假设。

硬件处理器在生产后基本不存在功能修复的可能。
所以这就要求在处理器设计仿真阶段能够进行功能的验证，也称为 “硅前”（Pre-silicon）验证。
现在常用的处理器验证手段有动态模糊测试与静态形式化验证两种。
这两种方法处理器验证手段也被应用到处理器侧信道漏洞检测上。

但是这些漏洞检测手段逐渐暴露出显著的局限性。
现代处理器的微架构状态空间极其庞大且交互复杂，基于动态仿真的模糊测试验证方法受限于有限的测试覆盖率，难以触发特定微架构状态下的深层硬件漏洞；
形式化方法凭借数学保证，能够基于寄存器传输级（RTL）或抽象硬件模型，进行严格的安全属性验证与详尽的状态空间探索。
理论上形式化方法能够穷尽所有侧信道漏洞。
但是状态爆炸，低效率，以及建模敏感成为了限制其功能发挥的关键。
如何让形式化检测框架在处理器流片前高效率地准确定位微架构缺陷仍是亟待研究的课题。
这正是本文提出“处理器侧信道漏洞形式化检测框架”的核心动机。


\subsection{本研究的意义和目的}

处理器特权分级机制是现代计算机系统的安全隔离的基石，但是处理器侧信道攻击打破了处理器这个对上层应用的安全保证。
处理器微架构中机密信息留下的痕迹可能通过指令执行时间、访存行为等侧信道被泄露出去。
同时，硬件设计漏洞难以在处理器流片后修复。
处理器的形式化验证方法可以在处理器设计时就找出这些漏洞，对处理器安全意义非凡。
但是随着现代处理器设计越发复杂，现有的处理器形式化验证方法面临着状态爆炸、建模不准确、安全属性通用性弱等问题。
导致处理器侧信道漏洞的形式化检测在现实应用中面临诸多问题。


这项工作旨在将验证对象扩展到处理器片上系统上（SoC），追求对真实世界处理器更真实的建模，检测出完整真实的处理器侧信道漏洞；
并且开发高效的同步状态注入技术，提高处理器侧信道漏洞形式化检测效率；
同时根据对新出现侧信道漏洞的洞察，提出通用性更强的处理器侧信道安全属性，扩展漏洞检测类型的覆盖面。
总的来说，作者希望该处理器侧信道漏洞形式化检测框架能够挖掘出更多触发链长、依赖微架构状态复杂、涉及部件多的深层次RISC-V处理器侧信道漏洞，同时达到更高的检测效率。